% !PRINTED 428
\weberpage{428}
\begin{center}
\S\ 118. \textbf{\gesperrt{Komposition der quadratischen Formen.}}
\end{center}

Da wir uns in der algebraischen Theorie der Klassengleichung auf die Theorie der Komposition der quadratischen Formen stützen müssen, so schicken wir darüber die folgenden Bemerkungen voraus:

Zwei primitive Formen $\varphi_1=(a_1,b_1,c_1)$, $\varphi_2=(a_2,b_2,c_2)$ der Diskriminante $-m$ heißen \gesperrt{einhellig} (einig, concordantes), wenn $a_1,a_2,\tfrac{1}{2}(b_1+b_2)$ ohne gemeinschaftlichen Teiler sind; $\varphi_1$ und $\varphi_2$ sind also gewiß einhellig, wenn schon $a_1,a_2$ ohne gemeinsamen Teiler sind, und da man, wenn nur die Klassen $k_1,k_2$ von $\varphi_1,\varphi_2$ gegeben sind, $\varphi_1,\varphi_2$ in ihren Klassen immer so wählen kann, so folgt, daß man in irgend zwei Klassen (die auch identisch sein können), immer zwei einhellige Formen finden kann. Wenn also diese Voraussetzung zutrifft, so folgt aus
\[
\tfrac{1}{4}(b_1^2-b_2^2) = \tfrac{1}{2}(b_1-b_2)\cdot\tfrac{1}{2}(b_1+b_2) = a_1c_1-a_2c_2,
\]
daß $\tfrac{1}{2}(b_1-b_2)$ durch den größten gemeinschaftlichen Teiler $\delta$ von $a_1,a_2$ teilbar ist, und daraus, daß man die beiden Kongruenzen
\weq{1}{b' \equiv b_1\ (\mathrm{mod}\ 2a_1), \quad b' \equiv b_2\ (\mathrm{mod}\ 2a_2)}
immer befriedigen kann.

Denn aus der ersten von ihnen folgt $b'=b_1+2\lambda a_1$, und man hat also $\lambda$ aus der Kongruenz
\[
\tfrac{1}{2}(b_1-b_2) + a_1\lambda \equiv 0\ (\mathrm{mod}\ a_2)
\]
zu bestimmen, die immer lösbar ist. Ist $\mu$ das kleinste gemeinschaftliche Vielfache von $a_1$ und $a_2$, also $a_1a_2=\mu\delta$, so kann $b'$ nach (1) noch durch $b=b'+2\mu h$ ersetzt werden, wenn $h$ eine beliebige ganze Zahl ist. Da nun $\tfrac{1}{4}(b_1^2+m)$ durch $a_1$, $\tfrac{1}{2}(b_2^2+m)$ durch $a_2$ teilbar ist, % [sic] printed with 1/4 and 1/2 respectively, inconsistently
so ist $\tfrac{1}{4}(b'^2+m)$ durch $\mu$ teilbar, und wenn wir $h$ aus der Kongruenz
\[
\frac{b'^2+m}{4\mu} + b'h \equiv 0\ (\mathrm{mod}\ \delta)
\]
bestimmen, die immer lösbar ist, da $b'$ nach Voraussetzung relativ prim zu $\delta$ ist, so wird $b^2+m=4a_1a_2c$ durch $4a_1a_2$ teilbar. Die zwei Klassen $k_1,k_2$ sind dann repräsentiert durch
\[
\varphi_1=(a_1,b,a_2c), \quad \varphi_2=(a_2,b,a_1c),
\]
und die Form
\[
\varphi = \varphi_1\varphi_2 = (a_1a_2,b,c)
\]
%CONT>>
